
This paper presents a measurement study of Expected Calibration Error (ECE) under adversarial conditions for an open-weight large language model (Llama-2-7b-hf). Evaluating on model-in-the-loop adversarial NLI benchmarks (ANLI R1--R3) and human-adversarial benchmarks (AdvGLUE MNLI, AdvGLUE QQP), this work finds that adversarial calibration degradation is not universal: it depends on how adversarial examples were constructed and on task type. AdvGLUE MNLI (human-adversarial NLI) produces a 7.1 percentage-point ECE increase relative to the clean MultiNLI baseline ($\Delta ECE$ = +0.071). ANLI R1 and R2 (model-in-the-loop adversarial NLI) produce calibration improvement ($\Delta ECE = -0.041$ and $-0.014$, respectively), while ANLI R3 produces modest degradation ($\Delta ECE$ = +0.024). Binary classification (AdvGLUE QQP) shows $\Delta ECE = -0.029$. Label preservation rates equal 1.000 for all adversarial splits by construction, confirming that $\Delta ECE$ reflects genuine calibration change and not label noise. RLHF alignment (Llama-2-7b-chat versus Llama-2-7b-hf base) consistently moderates calibration degradation on ANLI (100\% of rounds, $\Delta \Delta ECE$ up to +0.147) while exacerbating it on AdvGLUE MNLI ($\Delta \Delta ECE = -0.026)$. These findings, from a single-model pilot study, motivate adversarial calibration measurement as a distinct research problem from adversarial accuracy robustness, and provide a conditional framework---construction method $\times$ task type $\times$ RLHF alignment---for understanding when adversarial perturbation may degrade LLM confidence reliability.

